\documentclass[11pt]{article}

\usepackage[a4paper,margin=1in]{geometry}
\usepackage[T1]{fontenc}
\usepackage{lmodern}
\usepackage{graphicx}
\usepackage{amsmath,amssymb}
\usepackage{booktabs}
\usepackage{array}
\usepackage{tabularx}
\usepackage{longtable}
\usepackage{float}
\usepackage{subcaption}
\usepackage{hyperref}
\usepackage{enumitem}
\usepackage{xcolor}
\usepackage{microtype}
\usepackage{parskip}

\hypersetup{
    colorlinks=true,
    linkcolor=blue,
    urlcolor=blue,
    citecolor=blue
}

\setlength{\parskip}{6pt}

\title{%
  \textbf{Wafer Map Anomaly Detection on WM-811K}\\[0.4em]
  \large From Reconstruction Baselines to Pretrained Local Feature Modelling
}
\author{Deep Learning Project Report, Group 08\\[0.3em]
  Henry Lee Jun (1004219) \quad Chia Tang (1007200) \quad Genson Low (1005931)}
\date{\today}

\begin{document}

\begin{titlepage}
  \centering
  \vspace*{2cm}
  {\LARGE\bfseries Wafer Map Anomaly Detection on WM-811K\par}
  \vspace{0.5cm}
  {\large From Reconstruction Baselines to Pretrained Local Feature Modelling\par}
  \vspace{1.5cm}
  {\large\bfseries Deep Learning Project Report\par}
  \vspace{0.4cm}
  {\large 50.039 Deep Learning, Y2026\par}
  \vspace{0.4cm}
  {\large Group 08\par}
  \vspace{0.6cm}
  {\normalsize
    Henry Lee Jun (1004219) \quad
    Chia Tang (1007200) \quad
    Genson Low (1005931)\par}
  \vspace{1cm}
  {\large \today\par}
  \vfill
  {\small This report documents the experimental design, results, and diagnostic analysis
  for normal-only wafer anomaly detection on the WM-811K / LSWMD dataset.\par}
\end{titlepage}

\pagenumbering{roman}

\begin{abstract}
This report presents an anomaly-detection study on the WM-811K / LSWMD wafer-map dataset
under a strict normal-only training regime. The project begins with convolutional autoencoders
to establish a stable baseline pipeline, then moves progressively toward pretrained
ImageNet representations, teacher-student distillation, and PatchCore-style local anomaly
scoring. The main finding is that stronger frozen backbones combined with local feature
modelling improve performance substantially over plain reconstruction or global
embedding-distance baselines. The best completed run on the main benchmark is a
direct-$224\!\times\!224$ PatchCore model built on a pretrained ViT-B/16 backbone
(block 6 features, top-10\% patch reduction), achieving $\mathrm{F1}=0.595$,
$\mathrm{AUROC}=0.956$, and $\mathrm{AUPRC}=0.671$ under a deployment-style
95th-percentile validation threshold.
Score histograms and UMAP diagnostics show that the remaining bottleneck is
score-distribution overlap and naive thresholding rather than feature quality alone.
\end{abstract}

\newpage
\tableofcontents
\newpage

\pagenumbering{arabic}

\section{Introduction}

Semiconductor wafer fabrication spans hundreds of process steps, and a single
systematic error at any stage can render an entire region of dies unusable.
After fabrication, every die is electrically tested and the pass/fail result
recorded at its spatial location, producing a \emph{wafer bin map} (WBM): a
2-D spatial signature of the wafer's quality. Structured spatial patterns in
this map are diagnostically valuable, as a ring of failures near the wafer
edge implicates a different root cause than a central cluster or a diagonal
scratch. Automated recognition of such patterns enables rapid wafer disposition
and feeds back into upstream process control.

The na\"{i}ve approach of training a supervised classifier on labelled defect
examples faces a fundamental mismatch with real-world fab conditions. In the
WM-811K dataset used here, only $3.1\%$ of wafers carry a confirmed defect
label, $78.7\%$ are entirely unlabelled, and the labelled defect vocabulary
is not closed: novel failure modes emerge continuously as tools and processes
evolve. A supervised classifier trained on today's labelled patterns will be
blind to tomorrow's failures. These constraints motivate a \emph{normal-only
anomaly detection} framing: the model trains exclusively on normal wafers,
learns a representation of healthy wafer structure, and assigns an anomaly
score to each test wafer based on how far it deviates from that learned manifold.

This report documents a systematic progression through three modelling families,
all evaluated under the same normal-only training regime and deployment-style
threshold rule. Starting from convolutional autoencoder baselines, the work
moves through frozen pretrained backbone embeddings to local feature modelling
methods including teacher-student distillation, FastFlow, and PatchCore-style
nearest-neighbour patch scoring. The central finding is that neither a strong
backbone nor a strong local scoring method is sufficient alone; their combination,
at the right input resolution, drives the strongest results.

\section{Problem Framing and Manufacturing Context}

\subsection{Wafer Defects and Their Spatial Signatures}

During fabrication, a wafer passes through process steps including deposition,
lithography, etching, doping, and planarisation, repeated across dozens of
layers. Each step is a potential source of systematic error: equipment drift,
clean-room contamination, thermal gradients across the chuck, or lithographic
misalignment. When such an error occurs, the affected dies do not fail randomly;
they fail in spatially coherent clusters whose shape encodes the root cause.
A ring of failures near the wafer perimeter suggests edge-seal or etch
non-uniformity. A central cluster points to a chucking or thermal hot-spot
issue. A diagonal streak indicates mechanical contact from a probe or handler.
These spatial signatures are the basis of the wafer bin map and the information
this project aims to exploit automatically.

\subsection{Why Supervised Classification Falls Short}

The natural machine-learning response is supervised multi-class classification:
label each defect pattern type and train a classifier. This approach performs
well on clean, balanced benchmarks, but fails in practice for three reasons.

\textbf{Severe class imbalance.} Only $3.1\%$ of WM-811K wafers carry a
confirmed defect label. The remaining $96.9\%$ are either explicitly normal
($18.2\%$) or unlabelled ($78.7\%$, excluded from training to avoid label
noise). No classifier can operate reliably on this raw imbalance without
external rebalancing that distorts the real-world prevalence.

\textbf{Incomplete and evolving defect vocabulary.} The eight labelled pattern
types in WM-811K (Edge-Ring, Edge-Loc, Center, Loc, Scratch, Donut, Random,
Near-full) represent historically documented failures from a specific set of
fabs and process generations. Real production lines evolve continuously.
Novel failure modes emerge that no labelled example covers, and a supervised
classifier cannot flag what it has never seen.

\textbf{Extreme per-class label scarcity.} Even within the labelled defects,
some classes are critically sparse. The Scratch class provides only 15 test
examples and Near-full only 2. At this scale, per-class recall estimates are
unreliable and class-conditional training is impractical without synthetic data.

\subsection{Anomaly Detection as the Design Choice}

These constraints lead directly to a one-class \emph{anomaly detection}
framing. The model trains exclusively on wafers labelled as normal
(\texttt{failureType\,==\,none}) and learns the distribution of healthy
spatial structure. At inference time it assigns each wafer an anomaly score
reflecting how far its die-level pattern deviates from the learned normal
manifold. No defect labels are needed during training, the model is agnostic
to the defect vocabulary, and class imbalance is sidestepped by design.
The unlabelled majority of WM-811K, which may contain undetected defects,
is excluded from training entirely to keep the normal model uncontaminated.

\subsection{Operational Consequences of Decision Errors}

A deployed anomaly detector must commit to a binary decision for every wafer.
The two error types carry asymmetric consequences.

\begin{description}[leftmargin=2em, style=nextline]
  \item[False positive] A normal wafer is flagged as anomalous, triggering
        unnecessary manual review, reinspection, or scrap and reducing fab
        throughput. False positives are costly but recoverable.
  \item[False negative] A defective wafer passes inspection undetected,
        continuing through the production line and potentially being shipped
        to a customer. In safety-critical applications such as avionics or
        implantable medical devices, a missed defect can have catastrophic
        consequences. False negatives are generally the more serious failure mode.
\end{description}

Because a deployed system must commit to a specific operating threshold rather
than simply rank wafers, this report prioritises metrics that reflect real
threshold behaviour. A model with high AUROC but a poorly calibrated threshold
is operationally weak. The primary reported metric is therefore the F1 score
at a threshold derived from validation-normal scores alone, with no access to
defect labels at any point in threshold selection.

\section{Dataset and Split Construction}

\subsection{Source and Preprocessing}

The WM-811K dataset contains 811,457 wafer bin maps represented as variable-size
2-D matrices of die-level pass/fail values. The preprocessing pipeline retains
only explicitly labelled rows, converts each map to a normalised floating-point
array in $[0, 1]$, resizes using nearest-neighbour interpolation, and saves
the results alongside split metadata for reproducibility. Only wafers explicitly
labelled as normal are used for training; all other labelled failure types are
treated as anomaly. The unlabelled majority (${\sim}78.7\%$) is excluded
entirely to prevent label noise from entering evaluation.

\begin{table}[H]
\centering
\caption{WM-811K dataset composition.}
\label{tab:dataset_distribution}
\begin{tabular}{lcc}
\toprule
Category & Count & Fraction \\
\midrule
No-label & ${\approx}638{,}507$ & $78.7\%$ \\
Labelled: None (Normal) & $147{,}431$ & $18.2\%$ \\
Labelled: Pattern (Defect) & $25{,}519$ & $3.1\%$ \\
\midrule
\textbf{Total} & \textbf{811,457} & \textbf{100\%} \\
\bottomrule
\end{tabular}
\end{table}

Most experiments use wafer maps resized to $64\!\times\!64$ pixels. Later
pretrained-backbone experiments additionally evaluate direct $224\!\times\!224$
preprocessing to avoid the information loss that occurs when low-resolution
images are upsampled inside the model.

\subsection{Split Design}

All experiments share a single benchmark split drawing 50,000 normal wafers
from the labelled pool, with defect wafers held out exclusively in the test
partition. Normals are divided 80/10/10 across train, validation, and test;
defects are added only to the test set, capped at 5\% of the test-normal
count to reflect realistic production prevalence.

\begin{table}[H]
\centering
\caption{Main benchmark split.}
\label{tab:main_split}
\begin{tabular}{lcc}
\toprule
Subset & Normal wafers & Defect wafers \\
\midrule
Train   & 40{,}000 & 0   \\
Val     & 5{,}000  & 0   \\
Test    & 5{,}000  & 250 \\
\bottomrule
\end{tabular}
\end{table}

The validation split serves one purpose only: deriving the deployment
threshold. No defect labels are visible at any point before test evaluation.

\subsection{Defect Categories}

The 250 test-set defects span eight labelled pattern types that vary
substantially in spatial scale, coverage, and visual distinctiveness
(Table~\ref{tab:defect_types}, Figure~\ref{fig:defect_gallery}).

\begin{table}[H]
\centering
\caption{Defect-type breakdown of the main benchmark test set.}
\label{tab:defect_types}
\begin{tabularx}{\linewidth}{lc>{\raggedright\arraybackslash}X}
\toprule
Defect type & Count & Spatial character \\
\midrule
Edge-Ring & 84  & Continuous ring of failures near the wafer perimeter \\
Edge-Loc  & 53  & Localised cluster at or near the wafer edge \\
Center    & 50  & Compact cluster at the wafer centre \\
Loc       & 34  & Off-centre localised patch \\
Scratch   & 15  & Thin elongated streak across the die grid \\
Donut     & 7   & Curved partial ring inside the wafer body \\
Random    & 5   & Scattered failures with no spatial coherence \\
Near-full & 2   & Defective region covering nearly the entire wafer \\
\bottomrule
\end{tabularx}
\end{table}

\begin{figure}[H]
    \centering
    \includegraphics[width=\linewidth]{defect_type_gallery.png}
    \caption{Representative examples of all eight defect pattern types in the
    WM-811K evaluation set. Each image is a wafer bin map rendered as a
    binary die-level pass/fail grid. The diversity in spatial scale and
    structure, from full-wafer ring patterns to single-streak scratches,
    is a primary challenge for any global anomaly scoring approach.}
    \label{fig:defect_gallery}
\end{figure}

This diversity directly motivates the modelling choices described in
Section~\ref{sec:methodology}. Broad, globally visible defects such as
Edge-Ring and Center are detectable by models that measure overall
reconstruction quality. Small, localised defects such as Scratch and Loc
require spatially precise scoring that attends to individual patch regions
rather than image-level averages.

\section{Methodology}
\label{sec:methodology}

Table~\ref{tab:method_overview} summarises all model families evaluated in
this study. The methods are grouped by the research question they address and
progress from fully in-domain training to frozen pretrained features with
increasingly local scoring.

\begin{table}[H]
\centering
\caption{Overview of model families evaluated in this study.}
\label{tab:method_overview}
\begin{tabularx}{\linewidth}{>{\raggedright\arraybackslash}p{3.2cm}
                              >{\raggedright\arraybackslash}p{2.2cm}
                              >{\raggedright\arraybackslash}p{2.4cm}
                              >{\raggedright\arraybackslash}X}
\toprule
Method & Training & Scoring & Backbones / notes \\
\midrule
Convolutional Autoencoder
  & From scratch
  & Pixel reconstruction error
  & In-domain encoder; multiple architectural variants \\
Variational Autoencoder (VAE)
  & From scratch
  & Reconstruction error + KL divergence
  & KL weight $\beta$ swept \\
Deep SVDD
  & From scratch
  & Distance to latent hypersphere centre
  & One-class geometric objective; no reconstruction \\
\midrule
Pretrained backbone (center-distance)
  & None (frozen)
  & Global embedding distance to normal centroid
  & ResNet18; WideResNet50-2 \\
\midrule
Teacher-student distillation
  & Student trained on normals
  & Spatial discrepancy map (patch-level)
  & ResNet18; ResNet50; WideResNet50-2 \\
PatchCore
  & None (memory bank only)
  & Patch nearest-neighbour distance
  & ResNet18; ResNet50; WideResNet50-2; EfficientNet-B0 \\
FastFlow
  & Normalising flow trained on normals
  & Log-likelihood under learned feature distribution
  & WideResNet50-2 (multilayer) \\
\bottomrule
\end{tabularx}
\end{table}

\subsection{Evaluation Protocol}

Every model is evaluated under an identical protocol to keep comparisons fair
across model families. The primary operating threshold is derived from
validation-normal scores alone:
\[
  \tau_{0.95} = \text{95th percentile of validation-normal anomaly scores.}
\]
This threshold uses no defect labels and directly reflects a realistic
deployment constraint: in production, only normal wafers are available to
calibrate the detector. All reported precision, recall, and F1 values use
this threshold exclusively.

Five metrics are reported for each model. Val-threshold F1, the harmonic mean
of precision and recall at $\tau_{0.95}$, is the primary metric because it
measures performance at the actual deployed operating point. Precision and
recall at $\tau_{0.95}$ are reported separately to characterise whether a
model fails by flagging too many normal wafers (low precision) or by missing
too many defects (low recall). AUROC and AUPRC summarise ranking quality
across all thresholds and are particularly useful for comparing score
distributions, with AUPRC being the more informative of the two under
the class imbalance present in this dataset. A fifth metric, best-sweep F1,
reports the highest F1 achievable by sweeping thresholds on test labels.
This is an analysis ceiling only and is never used as a deployment result;
its value lies in quantifying how much headroom remains if a better threshold
selection rule were available.

\subsection{Modelling Approach and Motivation}

The experiments are organised around four progressively refined questions
about what drives anomaly detection performance on wafer maps. Full
training configurations and hyperparameter details for each method are
listed in Appendix~\ref{app:notebooks}.

\subsubsection{Question 1: Can reconstruction error detect wafer anomalies?}

The most direct anomaly detection strategy is to train a model to reconstruct
normal wafers and flag deviations at inference. If a model has only ever
seen normal patterns, the hypothesis is that it will reconstruct those
accurately while producing higher error on defective wafers whose patterns
were absent from training. Three model families test this idea.

\textbf{Convolutional Autoencoder.} An encoder-decoder trained end-to-end
on normal wafers to minimise pixel-wise reconstruction loss. The baseline
uses a latent dimension of 128, Adam optimisation with learning rate
$10^{-3}$ and weight decay $10^{-4}$, and early stopping on validation loss.
This is the natural first choice because it requires only in-domain data,
is fully trainable from scratch, and yields spatially resolved error maps.
A family of architectural variants tests whether the reconstruction ceiling
can be raised without changing the underlying paradigm: a batch-normalised
variant, a dropout sweep over rates $\{0.00, 0.05, 0.10, 0.20\}$, a residual
encoder-decoder with skip connections, and a higher-resolution
$128\!\times\!128$ version.

\textbf{Variational Autoencoder (VAE).} Extends the autoencoder by replacing
the deterministic latent code with a Gaussian latent distribution regularised
by a KL-divergence penalty. The KL weight $\beta$ is swept over
$\{0.001, 0.005, 0.01, 0.05\}$ to find the right trade-off between
reconstruction fidelity and latent regularisation. The motivation is that a
structured, probabilistic latent space may better capture the manifold of
normal wafers and produce a cleaner anomaly score.

\textbf{Deep SVDD.} Rather than using reconstruction error, Deep SVDD
trains a network to map all normal wafers into a compact hypersphere in
latent space. The anomaly score is the squared distance from a test
wafer's embedding to the learned hypersphere centre. This tests whether
a one-class geometric objective, with no reconstruction pathway at all,
can match or improve on the autoencoder family.

Across all three methods, a post-training scoring ablation compares seven
reduction rules applied to the same saved model checkpoint. These span global
pixel averages, foreground-masked variants that ignore background die positions,
a spatially smoothed pooled average, a single worst-pixel maximum, and a
top-$k$ percentile mean that averages only the highest-error pixels. This
ablation reveals that the scoring rule is a design decision co-equal with
architecture: the same trained model can rank very differently depending on
which score is applied.

\subsubsection{Question 2: Do pretrained ImageNet features help?}

Reconstruction models are limited by the capacity of a network trained
from scratch on a few tens of thousands of wafer images. A natural hypothesis
is that features from ImageNet pretraining encode richer, more transferable
visual structure and should produce better-separated embeddings for normal
versus anomalous wafers without any domain-specific training.

\textbf{Pretrained backbone center-distance baseline.} Frozen
ImageNet-pretrained backbones (ResNet18 and WideResNet50-2) are evaluated
using the simplest possible anomaly score: the L2 distance from each wafer's
global embedding to the centroid of train-normal embeddings. No fine-tuning
is applied. If pretrained features alone are sufficient, this baseline should
clearly outperform the trained autoencoders. The result deliberately stress-tests
the backbone-only hypothesis before any local scoring is introduced.

\subsubsection{Question 3: Does local spatial scoring matter?}

If global embedding distance is insufficient, the anomaly signal is likely
present in the spatial feature maps but diluted by global pooling. Small
localised defects occupy only a fraction of the wafer area; averaging their
contribution over the full image suppresses exactly the signal that matters.
Three methods test whether spatially precise, patch-level scoring recovers
that signal.

\textbf{Teacher-student distillation.} A frozen pretrained backbone (teacher)
extracts spatial feature maps from normal wafers. A lightweight student
convolutional network and an auxiliary feature autoencoder are trained jointly
to reproduce those feature maps on normal inputs only. At inference, the
per-location discrepancy between teacher and student forms a spatial anomaly
map: locations where the student fails to match the teacher are locally
anomalous. The map is reduced to a wafer score by averaging only the most
deviant patch locations, with the fraction of patches included swept as a
hyperparameter. Post-training branch weighting between the student and
autoencoder components is also swept to find the best deployed score
composition. Teacher backbones evaluated include ResNet18, ResNet50 (with
shallow and deep feature layers compared), and WideResNet50-2 (single-layer
and combined shallow-plus-deep).

\textbf{PatchCore.} Rather than training a student, PatchCore stores a
memory bank of patch-level feature embeddings extracted from normal wafers
during a single forward pass, then scores each test wafer by the nearest-neighbour
distance of its patch embeddings to the bank. No additional training is
required. Because each spatial patch is scored independently against its
closest normal neighbours, the method is inherently sensitive to localised
defects that affect only a small region of the image. Wafer-level scores
are reduced from per-patch distances using global mean, top-$k$ percentile
mean, or maximum patch distance, with memory bank size and the top-$k$ ratio
swept systematically. Backbones evaluated include ResNet18, ResNet50,
WideResNet50-2 with combined shallow and deep feature layers, and EfficientNet-B0.

\textbf{FastFlow.} An alternative local density estimator that fits a
normalising flow on top of frozen backbone feature maps. Rather than measuring
distance to stored neighbours, FastFlow assigns anomaly scores through
log-likelihood under the learned normal feature distribution. This tests
whether a generative density model over features can compete with the
nearest-neighbour approach of PatchCore without requiring a memory bank.
The WideResNet50-2 backbone with combined shallow and deep spatial feature
layers and four coupling-flow steps is the evaluated configuration.

\subsubsection{Question 4: Does input resolution matter for pretrained methods?}

Pretrained backbones expect $224\!\times\!224$ inputs. In the main
$64\!\times\!64$ experiments, wafer maps are cached at low resolution and
upsampled internally before entering the backbone, discarding spatial
information that was never captured. The final question is whether caching
wafers directly at $224\!\times\!224$ and feeding them natively to the backbone
provides a meaningful lift. This is evaluated for both the EfficientNet-B0
and WideResNet50-2 PatchCore branches as a controlled comparison against
their $64\!\times\!64$ counterparts, with all other hyperparameters held fixed.

\subsection{Experiment Staging}

The experiments follow a deliberate investigative sequence rather than a
parallel grid search. Each stage is designed to isolate one variable and let
its result motivate the next direction.

\paragraph{Stage 1: Reconstruction baseline.}
The plain autoencoder establishes the end-to-end pipeline, the validation-derived
threshold rule, and a first honest characterisation of the task difficulty.
Broad defects such as Edge-Ring and Center prove readily detectable; small
local defects such as Scratch and Loc expose the fundamental limit of
full-image reconstruction error as an anomaly signal.

\paragraph{Stage 2: Architecture and scoring ablations.}
BatchNorm, dropout, and residual variants, combined with the post-training
score sweep, determine whether the reconstruction ceiling can be raised
through in-domain tuning alone. The key finding from this stage is that
scoring rule choice is as impactful as architecture: the batch-normalised
checkpoint scored using maximum absolute pixel error outperforms the same
checkpoint scored using a top-$k$ mean, and reaches the highest
reconstruction-family F1. Dropout and higher resolution both fail to improve results.

\paragraph{Stage 3: Pretrained backbone baselines.}
The plain center-distance baselines for ResNet18 and WideResNet50-2 answer
Question 2 decisively. Both are weaker than the trained autoencoders despite
using much larger and richer feature spaces. The problem is not the backbone
quality but the global aggregation that collapses all spatial structure before
scoring.

\paragraph{Stage 4: Local anomaly scoring.}
Teacher-student and PatchCore experiments, staged from lighter to stronger
backbones, confirm that local spatial scoring is the decisive factor. Within
this stage, backbone scale (ResNet18 to ResNet50 to WideResNet50-2) provides
incremental gains, the choice of local aggregation ratio matters significantly,
and moving to direct $224\!\times\!224$ preprocessing provides a further clear
lift for the strongest backbones.

\paragraph{Stage 5: Diagnostics.}
Once the leading models are identified by metrics, UMAP embedding visualisations
and score histogram analysis examine why the fixed $\tau_{0.95}$ threshold
leaves headroom relative to the best-sweep F1. This diagnostic stage does not
change any model or result; it characterises the remaining bottleneck as
score-distribution overlap rather than feature quality.

\section{Main Results}

\subsection{Overall Leaderboard}

\begin{table}[H]
\centering
\caption{Selected top runs on the main benchmark split (ranked by val-threshold F1).}
\label{tab:main_leaderboard}
\begin{tabularx}{\linewidth}{>{\raggedright\arraybackslash}Xcccc}
\toprule
Model / variant & F1 & AUROC & AUPRC & Sweep F1 \\
\midrule
PatchCore-ViT-B16-x224-topk10-mb400k       & 0.595 & 0.956 & 0.671 & 0.650 \\
PatchCore-WideRes50-x224-topk-mb600k-r005 & 0.549 & 0.931 & 0.659 & 0.634 \\
PatchCore-EffNetB0-x224-topk-mb240k-r002  & 0.544 & 0.925 & 0.483 & 0.567 \\
PatchCore-WideRes50-topk-mb50k-r010       & 0.532 & 0.917 & 0.562 & 0.586 \\
Ensemble-PatchCore-TSRes50-maxpair        & 0.529 & 0.916 & 0.611 & 0.620 \\
TS-Res50-mixed-topk20                     & 0.525 & 0.909 & 0.599 & 0.606 \\
TS-WideRes50-multilayer-mixed-topk15      & 0.524 & 0.923 & 0.546 & 0.561 \\
FastFlow-WideRes50-l23-s4-mean            & 0.515 & 0.880 & 0.533 & 0.515 \\
AE-64-BN-max                              & 0.502 & 0.834 & 0.568 & 0.630 \\
\bottomrule
\end{tabularx}
\end{table}

\begin{figure}[H]
    \centering
    \includegraphics[width=1\linewidth]{images/autoencoder/overall_experiment_comparison.png}
    \caption{Overall comparison across completed experiments. Left: top runs ranked by deployed F1. Right: all completed runs in AUPRC-vs-F1 space, point size scaled by AUROC.}
    \label{fig:overall_comparison}
\end{figure}


\subsection{Three Cross-Cutting Findings}

\textbf{1. Score design matters almost as much as model design.}
The BatchNorm autoencoder with the default \texttt{topk\_abs\_mean} score ranked
below the plain baseline autoencoder ($\mathrm{F1}=0.431$ vs $0.468$).
Post-training rescoring with \texttt{max\_abs} lifted it to $\mathrm{F1}=0.502$,
making it the strongest reconstruction-family result without any retraining.
The same effect was observed in the teacher-student family: the imported
ResNet50 Kaggle checkpoint at its default score reached $\mathrm{F1}=0.488$,
while a post-training score sweep lifted it to $\mathrm{F1}=0.525$.

\textbf{2. Pretrained backbones need local anomaly scoring.}
Both the ResNet18 and WideResNet50-2 center-distance baselines were weak
($\mathrm{F1}<0.25$, $\mathrm{AUROC}<0.70$).  Applying PatchCore or
teacher-student distillation to the same backbone families produced the
strongest models in the report.

\textbf{3. Direct $224\!\times\!224$ preprocessing matters for pretrained
local methods.}
EfficientNet-B0 PatchCore rose from $\mathrm{F1}=0.467$ to $0.544$ when
wafer images were cached at native $224\!\times\!224$ instead of being
upsampled from $64\!\times\!64$.  WideResNet50-2 PatchCore similarly improved
over its $64\!\times\!64$ counterpart.

\subsection{Autoencoder Family}

\begin{figure}[H]
\centering
\includegraphics[width=\linewidth]{images/autoencoder/autoencoder_family_comparison.png}
\caption{Autoencoder family comparison. Left: all AE-family runs ranked by deployed F1.
Right: precision-recall operating points, point size scaled by AUPRC.}
\label{fig:ae_family}
\end{figure}

\begin{figure}[H]
\centering
\begin{subfigure}[t]{0.49\linewidth}
  \includegraphics[width=\linewidth]{images/autoencoder/ae64_training_and_evaluation.png}
  \caption{Baseline AE training curves and evaluation.}
\end{subfigure}
\hfill
\begin{subfigure}[t]{0.49\linewidth}
  \includegraphics[width=\linewidth]{images/autoencoder/ae64_score_ablation.png}
  \caption{Score ablation on the baseline AE checkpoint.}
\end{subfigure}
\caption{Baseline autoencoder ($64\!\times\!64$, latent dim 128): training behaviour
and the effect of scoring-rule choice on deployed F1 and AUROC.}
\label{fig:ae_training_ablation}
\end{figure}

The plain baseline autoencoder, trained for 43 epochs and scored with
\texttt{topk\_abs\_mean}, achieves $\mathrm{F1}=0.468$, $\mathrm{AUROC}=0.839$,
$\mathrm{AUPRC}=0.522$.  Key per-class results show that broad defects
(Edge-Ring recall 0.810, Center 0.720) are much easier than local ones
(Loc 0.206, Scratch 0.200).

Adding BatchNorm with the same scoring rule actually \emph{reduced} F1 to 0.431.
Score ablation on the BatchNorm checkpoint revealed that \texttt{max\_abs}
is a better match for the changed error distribution: rescoring lifts
F1 to 0.502, recall to 0.668, and AUPRC to 0.568, the best
reconstruction-family result in the report.

The dropout sweep ($0.00, 0.05, 0.10, 0.20$) confirmed that no dropout level
outperforms the no-dropout BatchNorm run.  Extending training epochs produced
only marginal changes, suggesting undertraining is not the main bottleneck.
Increasing resolution to $128\!\times\!128$ was slower and weaker.

\begin{figure}[H]
\centering
\includegraphics[width=1\linewidth]{images/autoencoder/autoencoder_resolution_comparison.png}
\caption{Resolution comparison: $64\!\times\!64$ vs $128\!\times\!128$ autoencoder.
Higher resolution did not improve anomaly detection on this dataset.}
\label{fig:ae_resolution}
\end{figure}

\subsection{VAE and Deep SVDD}

\begin{figure}[H]
\centering
\begin{subfigure}[t]{0.49\linewidth}
  \includegraphics[width=\linewidth]{images/vae/vae_beta_sweep.png}
  \caption{VAE beta sweep metrics.}
\end{subfigure}
\hfill
\begin{subfigure}[t]{0.49\linewidth}
  \includegraphics[width=\linewidth]{images/vae/svdd_training_and_evaluation.png}
  \caption{Deep SVDD training curves and evaluation.}
\end{subfigure}
\caption{Alternative one-class models. VAE beta tuning and Deep SVDD evaluation.}
\label{fig:vae_svdd}
\end{figure}

\begin{table}[H]
\centering
\caption{VAE and Deep SVDD results on the main benchmark split.}
\label{tab:vae_svdd}
\begin{tabular}{lcccc}
\toprule
Model & F1 & AUROC & AUPRC & Sweep F1 \\
\midrule
VAE ($\beta=0.005$) & 0.340 & 0.771 & 0.372 & 0.420 \\
VAE ($\beta=0.01$)  & 0.335 & 0.766 & 0.369 & 0.417 \\
Deep SVDD           & 0.360 & 0.788 & 0.213 & 0.366 \\
\bottomrule
\end{tabular}
\end{table}

The best VAE ($\beta=0.005$) improved over the initial $\beta=0.01$ run, but
remained clearly below the plain autoencoder.  Deep SVDD outperformed the
tuned VAE on F1 and AUROC but had very weak AUPRC (0.213), indicating
poor ranking quality under class imbalance.  Both models are useful
comparison references but not competitive for deployment.

\subsection{Pretrained Backbone Baselines}

\begin{figure}[H]
\centering
\includegraphics[width=\linewidth]{images/pretrained/compact_baseline_comparison.png}
\caption{Compact comparison of non-leading baselines: VAE, SVDD, and plain
embedding-distance backbones.}
\label{fig:compact_baselines}
\end{figure}

Both the frozen ResNet18 ($\mathrm{F1}=0.236$, $\mathrm{AUROC}=0.685$) and
WideResNet50-2 ($\mathrm{F1}=0.243$, $\mathrm{AUROC}=0.677$) center-distance
baselines confirm that a strong pretrained backbone alone cannot compensate
for globally aggregated scoring.

\subsection{PatchCore Family}

\begin{figure}[H]
\centering
\includegraphics[width=\linewidth]{images/patchcore/patchcore_family_comparison.png}
\caption{PatchCore family comparison across backbones and operating variants.
Left: best result per backbone. Right: all PatchCore runs coloured by
backbone/source and marked by wafer-level reduction.}
\label{fig:patchcore_family}
\end{figure}

PatchCore improved in stages.  The in-domain AE-BN-backed variant was weak
(best $\mathrm{F1}=0.336$, limited by the in-domain encoder's embedding quality).
Switching to ImageNet-pretrained ResNet18 produced a clear jump
($\mathrm{F1}=0.401$), and ResNet50 improved further ($\mathrm{F1}=0.420$),
showing the backbone direction was working.  Scaling to the multilayer
WideResNet50-2 ($\texttt{layer2+layer3}$) at $64\!\times\!64$ reached
$\mathrm{F1}=0.532$, the first PatchCore result to beat both the
AE family and the teacher-student baselines simultaneously.

The decisive finding about preprocessing: moving from the $64\!\times\!64$ cache
path to direct $224\!\times\!224$ preprocessing lifted the EfficientNet-B0
PatchCore from $\mathrm{F1}=0.467$ to $0.544$, and the WideResNet50-2 PatchCore
to $\mathrm{F1}=0.549$. A later ViT-B/16 direct-$224\!\times\!224$ PatchCore
follow-up pushed the main benchmark higher still, reaching $\mathrm{F1}=0.595$,
$\mathrm{AUROC}=0.956$, and $\mathrm{AUPRC}=0.671$.

Within each PatchCore sweep, \texttt{topk\_mean} reduction clearly outperformed
\texttt{mean} and \texttt{max} for the WideResNet50-2 backbone, and the best
top-k ratio band was consistently narrow around $r=0.05$--$0.10$.
Memory bank size mattered consistently: $50\,000$ (and the larger $600\,000$
for the $x224$ run) outperformed $10\,000$--$20\,000$.

\begin{table}[H]
\centering
\caption{WideResNet50-2 multilayer PatchCore sweep at direct $224\!\times\!224$.}
\label{tab:wrn_x224_sweep}
\begin{tabular}{lcccc}
\toprule
Variant & F1 & AUROC & AUPRC & Sweep F1 \\
\midrule
\texttt{topk\_mb600k\_r005\_x224} (selected) & \textbf{0.549} & \textbf{0.931} & \textbf{0.659} & \textbf{0.634} \\
\texttt{topk\_mb600k\_r010\_x224}            & 0.537 & 0.921 & 0.620 & 0.599 \\
\bottomrule
\end{tabular}
\end{table}

\begin{figure}[H]
\centering
\includegraphics[width=\linewidth]{images/patchcore/patchcore_sweep_barplots.png}
\caption{PatchCore sweep for WideResNet50-2 multilayer at direct $224\!\times\!224$
(validation-threshold F1 and AUROC for the two evaluated top-k ratio variants).}
\label{fig:wrn_x224_sweep}
\end{figure}

Per-defect recall for the saved WideResNet50-2 x224 diagnostic artifact:

\begin{table}[H]
\centering
\caption{Per-defect recall for the saved WideResNet50-2 x224 diagnostic artifact (r=0.05).}
\label{tab:wrn_x224_perdefect}
\begin{tabular}{lc}
\toprule
Defect type & Recall \\
\midrule
Donut     & 1.000 \\
Random    & 1.000 \\
Near-full & 1.000 \\
Loc       & 0.735 \\
Center    & 0.780 \\
Edge-Ring & 0.798 \\
Scratch   & 0.667 \\
Edge-Loc  & 0.623 \\
\bottomrule
\end{tabular}
\end{table}

\begin{figure}[H]
\centering
\begin{subfigure}[t]{0.55\linewidth}
  \includegraphics[width=\linewidth]{images/patchcore/18A2-patchcore-wideresnet50-multilayer-umap/score_histogram.png}
  \caption{Score histograms for the saved diagnostic artifact.}
\end{subfigure}
\hfill
\begin{subfigure}[t]{0.43\linewidth}
  \includegraphics[width=\linewidth]{images/patchcore/18A2-patchcore-wideresnet50-multilayer-umap/threshold_sweep_metrics.png}
  \caption{Threshold sweep: precision, recall, and F1 curves.}
\end{subfigure}
\caption{Score distribution and threshold sweep for the saved WideResNet50-2 PatchCore
x224 diagnostic artifact (\texttt{topk\_mb50k\_r005\_x224}). The 95th-percentile validation threshold
($\tau=0.537$) and the best-sweep threshold ($0.559$) are both marked.}
\label{fig:wrn_x224_diag}
\end{figure}

\subsection{Teacher-Student Distillation Family}

\begin{figure}[H]
\centering
\includegraphics[width=\linewidth]{images/teacherstudent/ts_family_comparison.png}
\caption{Teacher-student family comparison: ResNet18 vs ResNet50 teachers.
Left: deployed F1, AUPRC, and AUROC. Right: precision-recall operating points.}
\label{fig:ts_family}
\end{figure}

The teacher-student direction required careful score selection before becoming
competitive.  The ResNet18 teacher's default combined score was weak; rescoring
with student-only \texttt{topk\_mean} at ratio 0.20 lifted F1 to 0.495.  The
imported Kaggle-trained ResNet50 teacher at default score reached F1\,=\,0.488
with AUROC\,=\,0.913; local post-training score sweep (mixed branch weight
$2{:}1$, \texttt{topk\_mean}, ratio 0.20) lifted it to F1\,=\,0.525,
AUROC\,=\,0.909, AUPRC\,=\,0.599.

The WideResNet50-2 teacher produced the best teacher-student AUROC in the report
(multilayer $\texttt{layer2+layer3}$, mixed weights, ratio 0.15:
$\mathrm{AUROC}=0.923$, $\mathrm{F1}=0.524$).  The multilayer variant
outperformed the single-layer run consistently, confirming that combining two
teacher feature scales helps.

\begin{table}[H]
\centering
\caption{Teacher-student leaderboard entries on the main benchmark.}
\label{tab:ts_leaderboard}
\begin{tabular}{lccc}
\toprule
Variant & F1 & AUROC & AUPRC \\
\midrule
TS-Res50-mixed-topk20              & 0.525 & 0.909 & 0.599 \\
TS-WideRes50-multilayer-topk15     & 0.524 & 0.923 & 0.546 \\
TS-WideRes50-layer2-topk25         & 0.512 & 0.903 & 0.512 \\
TS-Res50-layer1-topk10             & 0.505 & 0.873 & 0.528 \\
TS-Res18-student-topk20            & 0.495 & 0.894 & 0.519 \\
\bottomrule
\end{tabular}
\end{table}

\subsection{WideResNet50-2 Family Overview}

\begin{figure}[H]
\centering
\includegraphics[width=\linewidth]{images/wrn/wrn_family_comparison.png}
\caption{WideResNet50-2 branch: plain embedding baseline, single-layer and multilayer
teacher-student runs, and multilayer PatchCore. Left: deployed F1 per selected row.
Right: precision-recall operating points, point size scaled by best-sweep F1.}
\label{fig:wrn_family}
\end{figure}

The WideResNet50-2 story follows a clean progression: the plain center-distance
baseline was the weakest result in the project ($\mathrm{F1}=0.243$); adding
teacher-student distillation jumped to $\mathrm{F1}=0.503$--$0.524$; switching
to multilayer PatchCore on the same backbone reached $\mathrm{F1}=0.532$ (64\,px)
and then $0.549$ (direct 224\,px).

\subsection{Cross-Model Per-Defect Recall Comparison}

\begin{table}[H]
\centering
\caption{Per-defect recall for selected representative runs.}
\label{tab:per_defect}
\begin{tabular}{lccccc}
\toprule
Defect type
  & AE-BN-max
  & TS-Res50-topk20
  & WRN-center
  & WRN-PatchCore-x64-r010
  & FastFlow-WRN-mean \\
\midrule
Scratch   & 0.467 & 0.333 & 0.200 & 0.533 & 0.200 \\
Loc       & 0.294 & 0.559 & 0.176 & 0.559 & 0.559 \\
Edge-Loc  & 0.472 & 0.491 & 0.132 & 0.585 & 0.547 \\
Center    & \textbf{0.800} & 0.720 & 0.140 & 0.620 & 0.720 \\
Donut     & 0.571 & 0.714 & 0.286 & \textbf{1.000} & 0.714 \\
Edge-Ring & 0.893 & 0.929 & 0.464 & 0.917 & 0.798 \\
Random    & 0.800 & \textbf{1.000} & 0.600 & \textbf{1.000} & \textbf{1.000} \\
Near-full & \textbf{1.000} & \textbf{1.000} & 0.000 & \textbf{1.000} & \textbf{1.000} \\
\bottomrule
\end{tabular}
\end{table}

The AE-BN-max run is the strongest for Center (0.800), a globally visible
defect well-suited to reconstruction error.  WRN PatchCore (x64) leads on
Scratch (0.533), Edge-Loc (0.585), and Donut (1.000).  The plain WRN
center-distance baseline struggles across the board, especially on Center
and Near-full.  Scratch remains the hardest class for all models.

\section{Expanded Holdout Validation}

To test whether leading checkpoints remain stable on a much larger evaluation
pool, a disjoint secondary holdout was constructed:

\begin{table}[H]
\centering
\caption{Main benchmark vs.\ expanded holdout dataset sizes.}
\label{tab:holdout_sizes}
\begin{tabular}{lcc}
\toprule
Subset & Main benchmark & Expanded holdout \\
\midrule
Train normals & 40{,}000 & 40{,}000 \\
Val normals   &  5{,}000 &  5{,}000 \\
Test normals  &  5{,}000 & 70{,}000 \\
Test defects  &    250   &  3{,}500 \\
\bottomrule
\end{tabular}
\end{table}

The defect pool expands approximately 14-fold, making rare-class recall estimates
far more meaningful.  The same train/validation normals are fixed across reruns,
so the checkpoint choice and threshold policy are unchanged.

\begin{table}[H]
\centering
\caption{Defect-type scaling from benchmark to expanded holdout.}
\label{tab:holdout_defects}
\begin{tabular}{lcc}
\toprule
Defect type & Benchmark & Holdout \\
\midrule
Edge-Ring & 84  & 1{,}302 \\
Edge-Loc  & 53  & 739     \\
Center    & 50  & 603     \\
Loc       & 34  & 492     \\
Scratch   & 15  & 169     \\
Random    &  5  & 108     \\
Donut     &  7  &  71     \\
Near-full &  2  &  16     \\
\bottomrule
\end{tabular}
\end{table}

\begin{table}[H]
\centering
\caption{Top expanded-holdout results among the reevaluated checkpoints
(70\,k normal / 3.5\,k defect test set).}
\label{tab:holdout_leaderboard}
\begin{tabularx}{\linewidth}{>{\raggedright\arraybackslash}Xcccc}
\toprule
Model / variant & F1 & AUROC & AUPRC & Sweep F1 \\
\midrule
TS-ResNet50                      & 0.516 & 0.926 & 0.622 & 0.585 \\
AE-BN (\texttt{max\_abs})        & 0.507 & 0.858 & 0.638 & 0.663 \\
Residual AE (\texttt{max\_abs})  & 0.503 & 0.871 & 0.637 & 0.643 \\
AE-BN-Dropout-0.00               & 0.501 & 0.857 & 0.654 & 0.673 \\
TS-ResNet18                      & 0.495 & 0.903 & 0.545 & 0.514 \\
PatchCore-ResNet50-mean-mb50k    & 0.426 & 0.824 & 0.373 & 0.449 \\
\bottomrule
\end{tabularx}
\end{table}

TS-ResNet50 leads this holdout bundle on deployed F1 and AUROC, demonstrating
strong stability on the larger test pool.  The best autoencoder variants trail
only slightly on deployed F1 but match or exceed TS-ResNet50 on AUPRC and
optimistic best-sweep F1, confirming they remain genuinely competitive.
The holdout false-positive rates stay close to the expected ${\approx}5\%$
implied by 95th-percentile thresholds, even with 70\,000 test-normal wafers
a reassuring stability signal.

This holdout bundle does not yet include the newer WRN $224\!\times\!224$
PatchCore winner from the main report.  It should therefore be read as a
robustness check on the reevaluated checkpoints, not as a full replacement
leaderboard.

\begin{figure}[H]
\centering
\begin{subfigure}[t]{0.32\linewidth}
  \includegraphics[width=\linewidth]{images/holdout70k_3p5k_evaluations/ts_resnet50.png}
  \caption{TS-ResNet50}
\end{subfigure}
\hfill
\begin{subfigure}[t]{0.32\linewidth}
  \includegraphics[width=\linewidth]{images/holdout70k_3p5k_evaluations/autoencoder_batchnorm__max_abs.png}
  \caption{AE-BN \texttt{max\_abs}}
\end{subfigure}
\hfill
\begin{subfigure}[t]{0.32\linewidth}
  \includegraphics[width=\linewidth]{images/holdout70k_3p5k_evaluations/patchcore_resnet50__mean_mb50k.png}
  \caption{PatchCore-ResNet50 mean}
\end{subfigure}
\caption{Score histograms from the expanded holdout (70\,k normal / 3.5\,k defect).
TS-ResNet50 shows the strongest deployed operating point among the reevaluated
checkpoints; the best autoencoder remains close; the best PatchCore rerun shows
heavier overlap near the threshold.}
\label{fig:holdout_histograms}
\end{figure}

\section{UMAP-Based Embedding Analysis}

UMAP projections are used as a geometric diagnostic after the leading models
are already identified by metrics.  The goal is to reveal the overlap structure
behind the score distributions and to explain why the fixed 95th-percentile
threshold is serviceable but still naive.
In this report, UMAP should be read primarily as a \emph{map of neighbourhood
structure}, not as a calibrated classifier by itself: points that land close
together have similar embeddings, while points that separate into different
branches or islands are geometrically less similar.  The split-coloured plots
show whether validation normals, test normals, and anomalies occupy the same
regions or drift into different parts of the manifold.  The score-coloured
plots show where the model assigns higher anomaly evidence inside that same
space.  The absolute axis values have no physical meaning; the useful signal is
the relative clustering, overlap, branch structure, and concentration of
high-score regions.  Where relevant, the same reference-fit UMAP is also used
to derive an auxiliary UMAP-space KNN score, calibrated on validation normals
only, as a geometry-aware alternative to a fixed score threshold.

\subsection{Plain WideResNet50-2 Embedding Baseline}

\begin{figure}[H]
\centering
\includegraphics[width=0.78\linewidth]{images/wrn/embedding_umap.png}
\caption{UMAP of the plain WideResNet50-2 embedding baseline. Anomaly wafers spread
through the same large islands as both validation and test normals.  A single
center-distance threshold is intrinsically limited when there is no geometric
separation.}
\label{fig:umap_wrn_baseline}
\end{figure}

The UMAP confirms the metric result: anomalies populate the dominant manifold
alongside normals, with no identifiable anomaly-only cluster.  A few anomaly-rich
side clumps exist, but the dominant picture is deep overlap, explaining why
center-distance scoring fails to separate classes.

\subsection{ResNet18 PatchCore}

\begin{figure}[H]
\centering
\begin{subfigure}[t]{0.48\linewidth}
  \includegraphics[width=\linewidth]{artifacts/x64/patchcore_resnet18_10A/max_mb50k/evaluation/plots/umap_by_split.png}
  \caption{Coloured by split}
\end{subfigure}
\hfill
\begin{subfigure}[t]{0.48\linewidth}
  \includegraphics[width=\linewidth]{artifacts/x64/patchcore_resnet18_10A/max_mb50k/evaluation/plots/umap_by_score.png}
  \caption{Coloured by anomaly score}
\end{subfigure}
\caption{Supplemental UMAP for the saved ResNet18 PatchCore \texttt{max\_mb50k}
artifact (qualitative geometry check; the selected best row is \texttt{mean\_mb50k}).
The reference-fit view is learned from sampled train-normal points and the
figure is exported as a capped subset ($5{,}000$ train reference, $4{,}000$
validation normals, $8{,}000$ test normals, $3{,}500$ test anomalies). Normals
and anomalies still share broad islands, but score-coloured heatmaps show
localised high-score pockets compatible with PatchCore local detection.}
\label{fig:umap_resnet18_patchcore}
\end{figure}

The split-coloured view shows that normals and anomalies continue to occupy
overlapping broad islands.  The score-coloured view, however, reveals localised
hotter regions inside the manifold: PatchCore has learned useful local patch
structure even where global geometry is mixed.  As an alternative thresholding
view, UMAP-space KNN was calibrated on validation normals only (95th percentile
threshold $=0.2337$), but this geometry-based score remained weak on the full
test set ($\mathrm{F1}=0.132$, $\mathrm{AUROC}=0.567$).  For ResNet18 PatchCore,
the UMAP is therefore best interpreted as a diagnostic visualisation rather than
as a stronger deployment threshold than the original PatchCore score.

\subsection{WideResNet50-2 PatchCore at Direct $224\!\times\!224$}

\begin{figure}[H]
\centering
\includegraphics[width=0.78\linewidth]{images/patchcore/18A2-patchcore-wideresnet50-multilayer-umap/umap_by_split.png}
\caption{UMAP for the selected WideResNet50-2 PatchCore $224\!\times\!224$
checkpoint (\texttt{topk\_mb50k\_r005\_x224}).  Validation and test normals
overlap heavily (reassuring: no distribution-shift artifact).  Anomalies now
concentrate more on branch edges, thinner tails, and smaller side islands
compared with the plain WRN embedding baseline, though full class separation
is not achieved.}
\label{fig:umap_wrn_patchcore}
\end{figure}

The geometry is healthier than the plain baseline: anomalies are no longer
uniformly distributed throughout the dominant manifold, but instead tend to
appear near structural boundaries and thinner branches.  This is compatible
with local patch scoring catching these out-of-manifold regions.  The remaining
overlap explains the moderate F1 values and the gap between 95th-percentile
and best-sweep thresholds.

\subsection{ViT-B/16 PatchCore at Direct $224\!\times\!224$}

\begin{figure}[H]
\centering
\begin{subfigure}[t]{0.48\linewidth}
  \includegraphics[width=\linewidth]{artifacts/x224/patchcore_vit_b16_5pct/main_5pct/umap_test_embeddings.png}
  \caption{Coloured by split}
\end{subfigure}
\hfill
\begin{subfigure}[t]{0.48\linewidth}
  \includegraphics[width=\linewidth]{artifacts/x224/patchcore_vit_b16_5pct/main_5pct/umap_by_score.png}
  \caption{Coloured by anomaly score}
\end{subfigure}
\caption{Reference-fit UMAP for the ViT-B/16 PatchCore main-benchmark run
(\texttt{vit\_b16\_one\_layer\_patchcore\_x224}). The manifold is fit on
$5{,}000$ sampled train-normal reference points and exported with the full
main-benchmark evaluation set ($5{,}000$ validation normals, $5{,}000$ test
normals, $250$ test anomalies).}
\label{fig:umap_vit_patchcore}
\end{figure}

The ViT geometry is useful as a diagnostic but not as a replacement threshold.
The split-coloured view shows that anomalies do not peel cleanly into one
isolated cluster; they still overlap heavily with the dominant normal manifold,
even though the score-coloured view highlights concentrated high-score regions
and sharper tails where the model assigns strong anomaly evidence.  This is
consistent with the main result: the original ViT PatchCore wafer score is very
strong on the main benchmark ($\mathrm{F1}=0.595$, $\mathrm{AUROC}=0.956$),
but a UMAP-space KNN threshold calibrated on validation normals remains weak
($\mathrm{F1}=0.083$, $\mathrm{AUROC}=0.525$).  For this branch, the UMAP is
therefore best interpreted as an explanatory embedding visualisation rather than
as a viable deployment thresholding method.  One implementation note is that
this ViT UMAP uses cosine distance, whereas the CNN PatchCore UMAPs in the
report use Euclidean distance.

\section{Thresholding Limitations and Discussion}

The 95th-percentile validation-normal threshold is fair, reproducible, and
requires no defect labels.  However, three consistent observations from the
diagnostics show its limitations:

\begin{itemize}[leftmargin=1.5em]
  \item \textbf{Score overlap.} Even the best models show substantial overlap
        between normal and anomaly score distributions.  The histogram for
        the WRN PatchCore x224 selected variant (Figure~\ref{fig:wrn_x224_diag})
        shows that the normal and anomaly distributions overlap near the
        operating threshold.
  \item \textbf{Geometric overlap.} UMAP figures show that anomalies are not
        cleanly peeled off from the normal manifold; they occupy boundaries
        and tails, but do not form an isolated cluster.
  \item \textbf{Sweep headroom.} Best-sweep F1 is systematically $0.06$--$0.08$
        above deployed-threshold F1 across the strongest models (e.g.\ WRN
        PatchCore x224: deployed 0.549, sweep 0.634).  This gap represents
        the threshold-selection headroom that a more principled rule could capture.
\end{itemize}

Possible directions to improve thresholding beyond the 95th-percentile rule:
\begin{itemize}[leftmargin=1.5em]
  \item Adaptive thresholding calibrated on validation score structure.
  \item Local-density or manifold-aware decision boundaries informed by embedding geometry.
        The current UMAP-space KNN follow-ups are not yet sufficient:
        for example, the ViT-B/16 branch keeps a very strong original wafer
        score while its 2-D UMAP-KNN score collapses to near-random ranking.
  \item Hybrid supervision: if any defect examples can be used for threshold tuning
        without contaminating training, they could inform a more calibrated operating point.
\end{itemize}

Not every main-report winner has been reevaluated on the expanded holdout yet.
The holdout bundle is therefore a stability check on the reevaluated subset,
not a universal replacement benchmark.

\section{Conclusion}

This project demonstrates a clear and reproducible progression in wafer-map
anomaly detection:

\begin{enumerate}[leftmargin=1.5em]
  \item Plain reconstruction models establish a viable pipeline and show that
        local-defect sensitivity is the main limitation.
  \item Better in-domain autoencoder variants and score selection improve
        performance without new training; the BatchNorm AE scored with
        \texttt{max\_abs} is the strongest reconstruction baseline
        ($\mathrm{F1}=0.502$).
  \item Pretrained backbones with global center-distance scoring are consistently
        weak regardless of backbone size.
  \item Local anomaly scoring methods (teacher-student distillation and
        PatchCore) produce the strongest results. The score composition
        choice interacts strongly with the checkpoint; post-training score
        sweeps are as important as architecture choices.
  \item Direct $224\!\times\!224$ preprocessing provides a clear lift for
        pretrained-backbone local methods.
  \item The best completed result on the main benchmark, ViT-B/16 PatchCore at
        direct $224\!\times\!224$, achieves $\mathrm{F1}=0.595$,
        $\mathrm{AUROC}=0.956$, and $\mathrm{AUPRC}=0.671$.
\end{enumerate}

The dual conclusion is: representation quality and local scoring matter
greatly, and the combination of a strong pretrained backbone with patch-level
nearest-neighbour scoring substantially outperforms reconstruction-based
or globally aggregated alternatives.  At the same time, the remaining
bottleneck is now score overlap and threshold selection rather than feature
extraction, which points toward more principled operating-point selection as
the highest-leverage next step.

\appendix
\newpage
\section{Appendix: Experiment Inventory}

This appendix documents the broader search space that sits behind the
streamlined main narrative.

\subsection{Autoencoder-Family Leaderboard (Main Benchmark)}

\begin{table}[H]
\centering
\caption{Selected autoencoder-family results (main benchmark, best score per checkpoint).}
\label{tab:appendix_ae}
\begin{tabular}{lccc}
\toprule
Variant & F1 & AUROC & AUPRC \\
\midrule
AE-64-BN-max              & 0.502 & 0.834 & 0.568 \\
AE-64-BN-DO0.00-max       & 0.487 & 0.851 & 0.617 \\
AE-64-BN-DO0.10-max       & 0.484 & 0.845 & 0.570 \\
AE-64-BN-DO0.05-max       & 0.483 & 0.835 & 0.552 \\
AE-64-Res-max             & 0.474 & 0.843 & 0.589 \\
AE-64-topk (43 ep)        & 0.460 & 0.835 & 0.525 \\
AE-64-BN-topk             & 0.431 & 0.790 & 0.603 \\
AE-128-mse                & 0.370 & 0.796 & 0.393 \\
\bottomrule
\end{tabular}
\end{table}

\subsection{PatchCore Leaderboard (Main Benchmark)}

\begin{table}[H]
\centering
\caption{Selected PatchCore results (main benchmark, best variant per backbone).}
\label{tab:appendix_patchcore}
\begin{tabular}{lccc}
\toprule
Variant & F1 & AUROC & AUPRC \\
\midrule
ViT-B16-x224-topk10-mb400k      & 0.595 & 0.956 & 0.671 \\
WideRes50-x224-topk-mb600k-r005 & 0.549 & 0.931 & 0.659 \\
EffNetB0-x224-topk-mb240k-r002  & 0.544 & 0.925 & 0.483 \\
WideRes50-topk-mb50k-r010       & 0.532 & 0.917 & 0.562 \\
WideRes50-topk-mb50k-r015       & 0.526 & 0.912 & 0.549 \\
WideRes50-topk-mb50k-r005       & 0.525 & 0.921 & 0.564 \\
EffNetB0-x64-topk-mb240k-r002   & 0.467 & 0.905 & 0.489 \\
Res50-mean-mb50k                & 0.420 & 0.821 & 0.363 \\
Res18-mean-mb50k                & 0.401 & 0.842 & 0.411 \\
AE-BN-mean-mb50k                & 0.336 & 0.851 & 0.226 \\
\bottomrule
\end{tabular}
\end{table}

\subsection{AE-Backed PatchCore Sweep}

\begin{table}[H]
\centering
\caption{PatchCore variants built on the frozen AE-BN encoder.}
\label{tab:appendix_aebn_patchcore}
\begin{tabularx}{\linewidth}{lcc>{\raggedright\arraybackslash}X}
\toprule
Variant & F1 & AUROC & Notes \\
\midrule
\texttt{mean\_mb50k}       & 0.336 & 0.851 & Best AE-BN PatchCore result \\
\texttt{topk\_mb50k\_r010} & 0.185 & 0.809 & Top-k reduction, 50\,k bank \\
\texttt{topk\_mb50k\_r005} & 0.123 & 0.777 & Narrower top-k region \\
\texttt{max\_mb50k}        & 0.056 & 0.679 & Max reduction too brittle \\
\bottomrule
\end{tabularx}
\end{table}

\subsection{ResNet18 Teacher-Student Focused Ablation}

\begin{table}[H]
\centering
\caption{ResNet18 teacher-student focused layer and top-k ablation.}
\label{tab:appendix_ts_resnet18}
\begin{tabular}{lcccc}
\toprule
Variant & Layer & Top-k ratio & F1 & AUROC \\
\midrule
\texttt{ts\_resnet18\_layer2\_topk25} & layer2 & 0.25 & 0.494 & 0.890 \\
\texttt{ts\_resnet18\_layer2\_topk20} & layer2 & 0.20 & 0.494 & 0.889 \\
\texttt{ts\_resnet18\_layer2\_topk15} & layer2 & 0.15 & 0.483 & 0.889 \\
\texttt{ts\_resnet18\_layer1\_topk15} & layer1 & 0.15 & 0.411 & 0.824 \\
\texttt{ts\_resnet18\_layer1\_topk20} & layer1 & 0.20 & 0.387 & 0.817 \\
\bottomrule
\end{tabular}
\end{table}

\subsection{ResNet50 Teacher-Student Score Sweep}

\begin{table}[H]
\centering
\caption{ResNet50 teacher-student local score-sweep summary (top rows).}
\label{tab:appendix_ts_resnet50}
\begin{tabular}{lcccc}
\toprule
Score variant & F1 & AUROC & AUPRC & Sweep F1 \\
\midrule
\texttt{s2\_a1\_topk\_mean\_r0.20} (selected) & 0.525 & 0.909 & 0.599 & 0.606 \\
\texttt{s1\_a0.5\_topk\_mean\_r0.20}           & 0.525 & 0.909 & 0.599 & 0.606 \\
\texttt{s1\_a2\_topk\_mean\_r0.20}             & 0.523 & 0.904 & 0.570 & 0.565 \\
\texttt{s2\_a1\_topk\_mean\_r0.10}             & 0.506 & 0.913 & 0.584 & 0.563 \\
\bottomrule
\end{tabular}
\end{table}

\subsection{WideResNet50-2 Teacher-Student Score Sweep}

\begin{table}[H]
\centering
\caption{WideResNet50-2 multilayer teacher-student score-sweep (top rows).}
\label{tab:appendix_ts_wrn}
\begin{tabular}{lcccc}
\toprule
Score variant & F1 & AUROC & AUPRC & Sweep F1 \\
\midrule
\texttt{s2\_a1\_topk\_mean\_r0.15} (selected) & 0.524 & 0.923 & 0.546 & 0.561 \\
\texttt{s1\_a1\_topk\_mean\_r0.25}             & 0.518 & 0.926 & 0.538 & 0.558 \\
\texttt{s2\_a1\_topk\_mean\_r0.20}             & 0.512 & 0.922 & 0.543 & 0.561 \\
\texttt{s2\_a1\_topk\_mean\_r0.10}             & 0.511 & 0.923 & 0.544 & 0.556 \\
\bottomrule
\end{tabular}
\end{table}

\subsection{FastFlow Ablation Summary}

\begin{table}[H]
\centering
\caption{FastFlow training-variant ablation (WideResNet50-2 backbone).}
\label{tab:appendix_fastflow}
\begin{tabular}{lccccc}
\toprule
Training variant & Layers & Steps & Best score & F1 & AUROC \\
\midrule
\texttt{wrn50\_l23\_s4} (selected) & layer2+layer3 & 4 & mean       & 0.515 & 0.880 \\
\texttt{wrn50\_l23\_s6}            & layer2+layer3 & 6 & mean       & 0.488 & 0.877 \\
\texttt{wrn50\_l2\_s6}             & layer2 only   & 6 & topk-r0.15 & 0.473 & 0.891 \\
\bottomrule
\end{tabular}
\end{table}

\subsection{Score Ensemble Study}

\begin{table}[H]
\centering
\caption{Post-hoc score ensemble: WRN PatchCore + TS-Res50.}
\label{tab:appendix_ensemble}
\begin{tabular}{lcccc}
\toprule
Fusion variant & F1 & AUROC & AUPRC & Sweep F1 \\
\midrule
PatchCore only         & 0.532 & 0.917 & 0.562 & 0.586 \\
max pair (selected)    & 0.529 & 0.916 & \textbf{0.611} & 0.620 \\
mean 70/30 PatchCore   & 0.528 & 0.918 & 0.596 & 0.616 \\
mean 50/50             & 0.526 & 0.917 & 0.606 & 0.615 \\
TS only                & 0.525 & 0.909 & 0.599 & 0.606 \\
\bottomrule
\end{tabular}
\end{table}

The ensemble improved AUPRC notably (0.611 vs 0.562) but did not beat the
standalone PatchCore on deployed F1.

\subsection{Experiment-to-File Reference}

\begin{table}[H]
\centering
\caption{Key experiment-to-notebook mapping for reproducibility.}
\label{tab:appendix_notebooks}
\begin{tabularx}{\linewidth}{>{\raggedright\arraybackslash}X>{\raggedright\arraybackslash}X}
\toprule
Experiment & Notebook / config \\
\midrule
Baseline AE                       & \texttt{02\_autoencoder\_training.ipynb} \\
BatchNorm AE                      & \texttt{05\_autoencoder\_batchnorm\_training.ipynb} \\
Dropout sweep                     & \texttt{06\_autoencoder\_batchnorm\_dropout\_training.ipynb} \\
Residual AE                       & \texttt{08\_autoencoder\_residual\_training.ipynb} \\
VAE beta sweep                    & \texttt{03\_vae\_training.ipynb} \\
Deep SVDD                         & \texttt{04\_svdd\_training.ipynb} \\
PatchCore AE-BN backbone          & \texttt{07\_patchcore\_training.ipynb} \\
ResNet18 backbone baseline        & \texttt{09\_resnet18\_backbone\_baseline.ipynb} \\
PatchCore ResNet18                & \texttt{10\_patchcore\_resnet18\_training.ipynb} \\
PatchCore ResNet50                & \texttt{11\_patchcore\_resnet50\_training.ipynb} \\
TS-ResNet18                       & \texttt{12\_ts\_distillation\_training.ipynb} \\
TS-ResNet50 (imported + analysis) & \texttt{13\_ts\_resnet50\_kaggle\_import\_analysis.ipynb} \\
WRN50-2 embedding baseline        & \texttt{15A\_wideresnet50\_2\_backbone\_baseline.ipynb} \\
WRN50-2 PatchCore x224            & \texttt{18A-patchcore-wideresnet50-multilayer-x224.ipynb} \\
WRN50-2 PatchCore x224 + UMAP    & \texttt{18A2-patchcore-wideresnet50-multilayer-x224\_with\_umap.ipynb} \\
EfficientNet-B0 PatchCore x64     & \texttt{21A\_patchcore\_efficientnet\_b0\_all-in-one.ipynb} \\
EfficientNet-B0 PatchCore x224    & \texttt{21B\_patchcore\_efficientnet\_b0\_all-in-one\_224.ipynb} \\
ViT-B/16 PatchCore x224           & \texttt{23\_patchcore-vit-b16-811k-one-layer.ipynb} \\
FastFlow sweep                    & \texttt{19A\_fastflow\_all\_in\_one.ipynb} \\
Score ensemble                    & \texttt{20\_score\_ensemble\_analysis.ipynb} \\
\bottomrule
\end{tabularx}
\end{table}

\end{document}
